% Generated draft: requires mathematical and source review.
\documentclass[11pt]{article}
\usepackage[margin=25mm]{geometry}
\usepackage{fontspec}
\setmainfont{Latin Modern Roman}
\usepackage{amsmath,amssymb,mathtools,unicode-math}
\setmathfont{Latin Modern Math}
\usepackage{enumitem,needspace,xurl,hyperref}
\setlength{\parindent}{0pt}
\setlength{\parskip}{5pt}
\setlength{\emergencystretch}{3em}
\newcommand{\N}{\mathbb{N}}
\newcommand{\Z}{\mathbb{Z}}
\newcommand{\Q}{\mathbb{Q}}
\newcommand{\R}{\mathbb{R}}
\newcommand{\C}{\mathbb{C}}
\newcommand{\F}{\mathbb{F}}
\newcommand{\A}{\mathbb{A}}
\newcommand{\PP}{\mathbb{P}}
\newcommand{\E}{\mathbb{E}}
\newcommand{\eps}{\varepsilon}
\newcommand{\dd}{\,\mathrm{d}}
\begin{document}
\section*{UnsolvedMath ID: 30002607}
\textbf{Source problem code:} catalogue entry
\subsection{Definitions and mathematical statement}
The supplied record concerns numerical homogenization: the construction of finite-dimensional approximation spaces for solutions of partial differential equations whose coefficients are rough or vary on multiple scales. A basis-design method selects basis functions \(\phi_1,\ldots,\phi_n\) and the associated space
\[
V_n=\operatorname{span}\{\phi_1,\ldots,\phi_n\}.
\]
The admissible function space, the meaning of roughness, and the admissible scale dependence are not specified in the record.

\textbf{Source question.} Can the identification of accurate basis elements for numerical homogenization of partial differential equations with rough or multiscale coefficients be facilitated by a general construction recipe or decision framework? In the supplied cleaned formulation, the question asks for a general method for constructing accurate finite-dimensional approximation bases for equations with arbitrary rough or multiscale coefficients.

To describe the missing accuracy specification, if a solution space \(X\), a norm \(\|\cdot\|_X\), and an exact solution \(u\in X\) were fixed, one possible measure of the approximation capacity of \(V_n\subseteq X\) would be
\[
\inf_{v\in V_n}\|u-v\|_X.
\]
This expression is explanatory notation only: the source does not select this error criterion, prescribe an error bound, or specify whether accuracy concerns best approximation or the output of a numerical scheme.

\textbf{Unresolved qualifications.} The source does not specify the spatial domain, equation class, coefficient assumptions, forcing data, initial or boundary conditions, solution concept, approximation norm, dimension-versus-error requirement, computational cost, or permitted dependence of the basis on the coefficients and data. Consequently, it does not determine a fully quantified existence statement or performance guarantee. The word ``arbitrary'' cannot be interpreted here as a specified class of well-posed equations. The original wording permits a decision framework, while the cleaned wording emphasizes a construction method; neither wording supplies a universal optimality criterion.
\subsection{Short English statement}
The problem asks for a general recipe or decision framework for designing accurate finite-dimensional bases for numerical homogenization with rough or multiscale PDE coefficients. The equation class, admissible coefficients, boundary conditions, and accuracy requirements remain unspecified, so the record supports a methodological question requiring further hypotheses before a precise quantified claim can be stated.
\subsection{Sources}
\begin{itemize}
\item Local UnsolvedMath registry: supplied complete record with permanent ID \(30002607\), title \emph{Basis Design for Numerical Homogenization}, exact source problem code \texttt{catalogue entry}, in the \emph{Oberwolfach Reports Open Problems} collection. The supplied proposed year is 2014; no individual proposer is identified.
\item Supplied report citation: \emph{Computational Multiscale Methods (2015)}, Oberwolfach Reports, pp.\ 1625--1681, DOI \texttt{10.4171/owr/2014/30}. \url{https://doi.org/10.4171/owr/2014/30}. The background attributes the problem to workshop open-problem discussions documented in Oberwolfach Reports. The citation's 2015 designation is retained alongside the supplied 2014 report identifier and proposed year; their relationship is not resolved in this draft.
\item Supplied literature reference: Owhadi, \emph{Bayesian numerical homogenization}. \url{https://arxiv.org/abs/1406.6668}.
\item Supplied literature reference: \emph{Oberwolfach Report 30/2014}. \url{https://ems.press/content/serial-article-files/46520}.
\item Supplied literature assessment, checked \textbf{2026-08-21}: the record labels the problem \emph{partially solved} and states that Bayesian numerical homogenization supplies a general decision framework and derives optimal-recovery bases for broad integro-differential problems, while a universally optimal basis-design recipe for arbitrary multiscale operators is not established. This is supplied dated curation, described in the background as literature triage derived from the supplied OWR curation; it has not been independently verified here. Its time-sensitive status and cited claims require independent checking.
\item The supplied statement status is \texttt{corrected\_verified}. Its verification note reports checking against the official Oberwolfach report and says that the cleaned statement restores the source's hypotheses and notation while preserving its open questions. This is a supplied provenance assertion, not independent verification performed for this draft.
\end{itemize}
\end{document}
